\section{Introduction}
\label{sec:intro}

Locating code that provides evidence for a written requirement is a recurring need in software analysis, for example when checking whether a repository supports the requirements of a task. 
Tracing requirements to code has a long history in software engineering~\cite{antoniol2002traceability,gao2024triad,kuang2023analyzing}. 
We study a traceability-style variant: given a technical requirement phrased in natural language, here a skill required by a job description (JD), retrieve the code chunks that evidence it.

Chunking via abstract syntax trees (ASTs) is motivated by the observation that line-based splitting can break semantic structures~\cite{zhang2025cast}. 
Recent controlled evidence in code completion, however, finds that structure-aware chunking does not outperform a simple sliding window~\cite{wu2026chunking}, and another study of code-RAG design choices reports that line-based chunking matches syntax-aware splitting across context budgets~\cite{galimzyanov2025practical}. 
Whether this holds when the goal is to retrieve evidence for a requirement is unclear to us. 
Moreover, real-world source files can contain constructs, such as declarative configurations, UI templates, and type declarations, that standard method-level AST extraction does not cover, and chunk boundaries are often confounded with the metadata (e.g., file path, class name) prepended to chunks. 
To the best of our knowledge, no prior study crosses chunking strategy with context headers for requirement-to-code evidence retrieval with human relevance labels.

In this Research Preview, we investigate how chunk boundary formulation and contextual metadata interact to determine retrieval effectiveness. 
We formulate four research questions:
\begin{itemize}[leftmargin=*,itemsep=1pt,topsep=1pt]
    \item \textbf{RQ1 (Preliminary):} Does prepending a hierarchical context header improve the retrieval of requirement evidence from source code chunks?
    \item \textbf{RQ2 (Preliminary):} Does AST-first chunking outperform line-window chunking when both conditions use, or both omit, context headers?
    \item \textbf{RQ3 (Planned):} Which context header components drive retrieval improvements, and how do chunking trade-offs generalize across programming paradigms?
    \item \textbf{RQ4 (Planned):} Does retrieved code evidence improve human evaluators' accuracy and confidence in downstream requirement verification?
\end{itemize}

To address RQ1 and RQ2, we run a controlled $2 \times 2$ factorial experiment on 25 public JDs and 40 open-source repositories (20 Java Spring Boot, 20 React/TypeScript; 798 chunks) at three retrieval tiers (BM25, dense BGE-M3, and cross-encoder reranking). 
Preliminarily, headers improve reranking of AST chunks ($\Delta = +0.112$, $p_{\mathrm{Holm}} = 0.0031$); without headers, line windows outperform AST chunks ($\Delta = -0.126$, $p_{\mathrm{Holm}} = 0.0001$); and we find no evidence that AST chunks with headers differ from line windows without headers ($\Delta = -0.014$, $p_{\mathrm{Holm}} = 0.85$).

We make four contributions:
\begin{enumerate}[leftmargin=*,itemsep=1pt,topsep=1pt]
    \item A factorial design that crosses chunking strategy with context headers for requirement-to-code evidence retrieval over multi-file repositories.
    \item Preliminary findings: headers improve reranking of AST chunks, headerless line windows outperform AST chunks, and we find no evidence that AST chunks with headers differ from headerless line windows; an audit shows that 35.7\% of AST-arm chunks are line-window fallbacks.
    \item A replication package with 603 human-labeled JD-chunk pairs (weighted kappa 0.78 overall, 0.88 on the initial gold subset, and 0.68 on the pooled items) and retrieval runs.
    \item A research plan to test why headers help and whether retrieved evidence supports human verification.
\end{enumerate}
The study relates to SANER topics on software analysis, mining software repositories, and program comprehension.
